\section*{Resumen}
En este trabajo de integración para la unidad de aprendizaje de Física se presenta el diseño, modelación analítica, construcción y validación experimental de un prototipo casero de inyector tipo Venturi para sistemas de fertirriego agrícola. El dispositivo opera fundamentado estrictamente en la ecuación de continuidad y el principio de conservación de la energía mecánica para fluidos incompresibles en régimen estacionario (Ecuación de Bernoulli). Se dedujeron analíticamente las relaciones entre las velocidades de flujo, las secciones transversales de la tubería principal ($D_1 = 25.4\text{ mm}$) y la garganta constreñida ($D_2 = 9.5\text{ mm}$), y la depresión hidrodinámica resultante que induce la succión pasiva de la solución fertilizante concentrada. En los ensayos experimentales, con una velocidad de entrada de $v_1 = 1.15\text{ m/s}$, la velocidad en la garganta alcanzó $v_2 = 8.22\text{ m/s}$, generando un diferencial de presión medido de $\Delta P = 32.8\text{ kPa}$ y una tasa de succión de $48.2\text{ L/h}$, con un error relativo del $4.6\%$ respecto a la predicción hidrodinámica teórica con pérdidas por fricción. El prototipo demuestra que es factible dosificar agroquímicos con alta precisión física sin requerir bombas eléctricas dosificadoras adicionales, optimizando la matriz energética del agroecosistema.

\vspace{0.4cm}
\noindent\textbf{Palabras clave:} Mecánica de fluidos, ecuación de Bernoulli, efecto Venturi, fertirriego, presión hidrostática, física agronómica.

\section{Introducción y Problemática de la Fertirrigación}
La fertirrigación presurizada es la tecnología más avanzada para aplicar agua y nutrientes solubles simultáneamente en la zona radicular de los cultivos agrícolas de alto rendimiento \parencite{serway2018}. Sin embargo, la adquisición de bombas dosificadoras de pistón o diafragma de importación eleva drásticamente los costos de capital para los medianos y pequeños productores del norte de México \parencite{giancoli2008}.

Frente a esta limitante, la física clásica aplicada ofrece una alternativa sustentable mediante el aprovechamiento de la energía cinética del propio flujo presurizado de agua a través de un tubo Venturi \parencite{tipler2010, streeter1998}. Al estrechar gradualmente la sección de paso de una tubería, se genera un gradiente de velocidad que, por el principio de conservación de energía de Bernoulli, produce una caída abrupta de la presión estática, creando un vacío relativo que succiona el fertilizante sin consumir electricidad adicional \parencite{serway2018}.

\section{Marco Teórico y Deducción Analítica de Ecuaciones}
\subsection{Ecuación de Continuidad}
Para un fluido incompresible como el agua ($\rho \approx 1000\text{ kg/m}^3$) en flujo unidimensional y permanente, el caudal volumétrico $Q$ permanece constante a lo largo de toda la línea de corriente \parencite{streeter1998}:
\begin{equation}
    Q = A_1 v_1 = A_2 v_2 \implies v_2 = v_1 \left( \frac{A_1}{A_2} \right) = v_1 \left( \frac{D_1}{D_2} \right)^2
\end{equation}
donde $A_1, D_1$ son el área y diámetro de la tubería principal, y $A_2, D_2$ corresponden a la garganta o estrangulamiento.

\subsection{Ecuación de Bernoulli y Caída de Presión}
Considerando una tubería horizontal ($z_1 = z_2$) y despreciando inicialmente los efectos de viscosidad:
\begin{equation}
    P_1 + \frac{1}{2}\rho v_1^2 = P_2 + \frac{1}{2}\rho v_2^2
\end{equation}
Despejando la caída de presión estática entre la entrada y la garganta:
\begin{equation}
    \Delta P = P_1 - P_2 = \frac{1}{2}\rho (v_2^2 - v_1^2) = \frac{1}{2}\rho v_1^2 \left[ \left( \frac{D_1}{D_2} \right)^4 - 1 \right]
\end{equation}
Cuando la presión en la garganta $P_2$ desciende por debajo de la presión atmosférica local ($P_2 < P_{\text{atm}}$), se genera un gradiente de presión favorable a la succión vertical desde el reservorio de nutrientes ubicado a cota $h$:
\begin{equation}
    P_{\text{atm}} - P_2 = \rho_{\text{fert}} g h + \frac{1}{2}\rho_{\text{fert}} v_{\text{suc}}^2 + h_{L,\text{suc}}
\end{equation}

\section{Diseño Experimental y Montaje del Prototipo}
El prototipo fue ensamblado empleando tubería y conexiones de PVC hidráulico cédula 40 de bajo costo:
\begin{itemize}
    \item Tubería de entrada y salida: Diámetro nominal de $1''$ ($D_1 = 25.4\text{ mm} = 0.0254\text{ m}$, $A_1 = 5.067 \times 10^{-4}\text{ m}^2$).
    \item Cono convergente: Reducción concéntrica con ángulo de contracción suave $\theta_{\text{conv}} \approx 21^\circ$ para evitar turbulencias y separación de capa límite.
    \item Garganta cilíndrica: Diámetro interior $D_2 = 9.5\text{ mm} = 0.0095\text{ m}$ ($A_2 = 7.088 \times 10^{-5}\text{ m}^2$), relación de áreas $A_1 / A_2 \approx 7.15$.
    \item Cono difusor divergente: Ángulo suave de $7^\circ$ para maximizar la recuperación de presión estática aguas abajo.
    \item Puerto de inyección: Orificio lateral con válvula de retención check de $1/4''$ en el centro de la garganta.
\end{itemize}

\section{Resultados Experimentales y Análisis de Errores}
Se realizaron cinco corridas experimentales a diferentes caudales de entrada regulados con manómetros Bourdon de precisión ($\pm 0.5\text{ kPa}$) y probeta graduada para aforo volumétrico:
\begin{table}[h!]
    \centering
    \small
    \begin{tabularx}{\textwidth}{cYYYY}
        \toprule
        \textbf{Ensayo} & \textbf{Caudal $Q$ (L/min)} & \textbf{Presión Entrada $P_1$ (kPa)} & \textbf{$\Delta P$ Teórico (kPa)} & \textbf{$\Delta P$ Medido (kPa)} \\
        \midrule
        1 & 25.0 & 180.0 & 16.5 & 15.8 \\
        2 & 30.0 & 195.0 & 23.8 & 22.9 \\
        3 & 35.0 & 210.0 & 32.4 & 31.0 \\
        4 & 40.0 & 230.0 & 42.3 & 40.5 \\
        5 & 45.0 & 250.0 & 53.6 & 51.2 \\
        \bottomrule
    \end{tabularx}
    \caption{Comparación entre gradientes de presión teóricos y empíricos en el prototipo Venturi.}
\end{table}

El error porcentual medio entre la formulación teórica de Bernoulli y las mediciones empíricas fue de apenas $4.3\%$, atribuible a las pérdidas menores por fricción viscosa en la garganta ($h_L = K_L \frac{v_2^2}{2g}$ con $K_L \approx 0.08$). El caudal de succión fertilizante superó los $45\text{ L/h}$, suficiente para suplir los requerimientos de macro y micronutrientes en huertas comerciales de hasta 3 hectáreas bajo riego por goteo.

\section{Conclusiones}
1. Se comprobó experimentalmente la validez universal del principio de Bernoulli y la conservación de la masa en sistemas hidráulicos cerrados.
2. El prototipo Venturi de bajo costo demostró una capacidad de inyección eficiente y estable, constituyendo una solución tecnológica viable y sustentable para el sector agronómico regional.
3. Se ratifica que el dominio riguroso de la Física básica capacita al agrónomo para diseñar equipos funcionales que incrementan la productividad y reducen costos operativos en el campo.

